\centerline{\bf \Huge ОТА. Делимость I}
\centerline{\bf \Huge Решения}

\begin{flushright}
        \scriptsize{Паника — признак слабости духа\\Пешши из ДжоДжо}
\end{flushright}

\problem{1}
Докажите, что произведение любых трёх последовательных натуральных чисел делится на 6.

\textbf{Решение}
Среди этих трёх чисел есть хотя бы одно чётное число. Значит, в разложении произведения на простые множители есть множитель 2. Среди этих трёх чисел одно число делится на 3 . Значит, в разложении произведения на простые множители есть множитель 3. В разложении произведения на простые множители есть простые числа 2 и 3 . Значит, оно делится на их произведение, то есть на 6.

\textit{Для ассистентов:} Попросите доказать, что один из множителей делится на 2 и один из множителей делится на 3.

\problem{2}
Про $a, b$ известно, что ни одно из них не кончается на ноль, а произведение 10000 . Какие значения может принимать их сумма?

\textbf{Решение.} Разложим 10000 на простые множители. Получится $10000 = 2^4\cdot5^4$. Любой делитель числа 10000 имеет вид $2^x\cdot5^y$. Но нам нужны делители не делящиеся на 10. Такие только $2^4$ и $5^4$. Поэтому их сумма может быть только $16+625 = 641$

\textit{Для ассистентов:} Попросите их разложить число 10000 на простые множители и спросите. Спросите, на что должно делиться число, чтобы оно оканчивалось на 10 и что должно быть в разложении на простые множители.

\problem{3}
Существуют ли 3 натуральных числа таких, что ни одно из них не делится на другое, произведение любых двух из них делится на третье?

\textbf{Решение.} Да существуют. Например, числа 6, 10, 15.

\problem{4}
Существует ли n такое, что $n!$, что нулей на конце ровно $5?$

\textbf{Решение.} 
4! оканчивается на четыре нуля. При добавлении множителя 25 добавится еще два нуля (так как число 25 в разложение факториала добавит сразу две пятерки). Поэтому 25! оканчивается уже на шесть нулей.

\textit{Для ассистентов:}
Объясните ученику за счёт чего появляются нули на конце у числа $n!$. Покажите на примере 10! или 15!

\problem{5}
Если в выражение $\mathrm{n}^2-\mathrm{n}+41$ подставлять числа $\mathrm{n}=1,2,3,4,5$, то получатся простые числа $41,43,47,53,61$. Верно ли, что при любом n получится простое число?

\textbf{Решение.} Нет, неверно. При $n=41$ получится составное число.

\textit{Для ассистентов:} ну хз тут реально отсталым надо быть, чтобы не решить

\problem{6}
Докажите, что у составного числа $a$ найдется такой простой делитель $p$, что $p^2 \leqslant a$.

\textbf{Решение.} Пусть его нет. Возьмём любой простой делитель $p > \sqrt{a}$. Тогда $\dfrac{a}{p}$ тоже делитель $a$ и тоже должен быть больше $\sqrt{a}$. Но тогда $p \cdot \dfrac{a}{p} > \sqrt{a} \cdot \sqrt{a} = a$. Противоречие.

\textit{Для ассистентов:} попросите доказать от противного